\setcounter{equation}{0}
\section{Parallel Simulation}

In this thesis, the LAM message passing interface (MPI) has been used distribute the computation (maintained by Indiana university)\footnote{http://www.lam-mpi.org/}. The program is SIMD (Single Instruction Multiple Data), with each node running a copy of the same code, but operating on different datas. The heaviest computational burden in the simulation of a quantum computer is the matrix multiplication when a register is applied a quantum program. The resulting rows of the output vector contain n multiplications, and n-1 additions for a quantum register with n qubits. Although all nodes will be running the same code, one of them will perform special tasks, it will be used to distribute the data (input vector and quantum program) and get the results (output vector). This node will be called the master node.
A strategy is to keep the probability amplitudes of each qubit on different processors, and apply the quantum gates to individual qubits. However this approach doesn't let us obtain the output vector after the execution of the quantum program, so it is required to collect the qubit's values in a single processor and perform a tensor product to get the output vector, which implies serious communication and computation overheads.
Another strategy is to assign each row of the input vector to individual nodes, and perform the computation of the corresponding row of the outcome vector on this node. Since only one row of the Quantum program is needed to calculate the output vector's row, the quantum program can be efficiently distributed. However, with this method, the nodes have to communicate the members of the input vector with each other, which increases the communication overhead, and code complexity.
Since all the nodes must have all the members of the input vector to calculate the corresponding row of the output vector, it is quite logical to distribute a copy of the input vector to all processors before execution. The program supplied with this thesis uses this strategy, it broadcasts the input vector to all nodes, to minimize the communication overhead. Then the quantum program is divided into chunks of appropriate sizes and distributed to corresponding nodes. With this strategy, it is possible to efficiently distribute the bulk of the computation among the worker nodes. After the computation is finished, the nodes begin broadcasting their results. Actually, if each node sent it's result to the master node, it would be possible to get the output vector on the master node. However, if the quantum programs are applied sequentially, this approach needs to send the output vector to the worker nodes as the input vector of the next quantum program, thus increasing the communication overhead. So it is wise to keep a copy of the output vector on all nodes to minimize the communication overhead. Depending on the network architecture, it may be more efficient to use a hypercube template to gather all the data on all the nodes (if nodes are connected with a hypercube topology). However, with a classical star topology, plain broadcasts of the results will be more efficient.
